\section{Contexto e Problema}
\label{sec:introducao}

Sistemas de estacionamento, pedágio, cadastro de veículos e
seguros recebem placas digitadas por pessoas. Erros de digitação,
caracteres trocados e formatos misturados são comuns, e cada falha
só aparece depois, quando a consulta ao cadastro não encontra o
veículo.

O Brasil convive com dois formatos de placa: o anterior ao
Mercosul, com três letras e quatro algarismos, e o Mercosul, com
três letras, um algarismo, uma letra e dois algarismos
\cite{contran780,contran741}. Um validador que confere o formato
antes da consulta evita requisições inúteis e dá retorno imediato
a quem digita.

\subsection{O que foi entregue}

\begin{itemize}
  \item uma especificação precisa do que é e do que não é uma placa
        válida;
  \item um modelo matemático (autômato finito determinístico) que
        decide a validade de forma exata e previsível;
  \item uma implementação em Python que segue o modelo e foi verificada
        com casos válidos e inválidos;
  \item a demonstração de que o modelo não tem estados desnecessários.
\end{itemize}

\subsection{O que não faz parte do escopo}

O validador reconhece a \emph{estrutura} da sequência de
caracteres. Elementos visuais da chapa (cor, faixa azul, QR Code,
bandeira, logotipo e tipografia) não são avaliados, e reconhecer o
formato não comprova que a placa esteja registrada ou vinculada a
um veículo.

\subsection{Como ler este documento}

A Seção~\ref{sec:como-funciona} explica a ideia da solução.
A Seção~\ref{sec:especificacao} define o que o validador aceita e
rejeita, com exemplos. As Seções~\ref{sec:modelagem} a
\ref{sec:minimizacao} trazem o modelo matemático, os diagramas e a
garantia de que o autômato é mínimo. A Seção~\ref{sec:testes}
mostra os testes, e a Seção~\ref{sec:limitacoes} discute
limitações e evolução.